\chapter{经典场论基础}

\section{从粒子力学到场论}

在有限自由度的粒子力学中，体系的动力学变量是有限个广义坐标
\[
q^{i}=q^{i}(t).
\]
这些广义坐标只依赖于时间。相应的作用量通常写为
\[
S=\int_{t_{1}}^{t_{2}}L(q^{i},\dot q^{i},t)\,\dif t.
\]
真实运动轨道由作用量驻值条件
\[
\delta S=0
\]
确定。

在经典场论中，动力学变量不再是有限个广义坐标，而是定义在时空每一点上的场变量。设场分量记为
\[
\phi^{a}=\phi^{a}(x^{\mu}),
\]
其中指标 $a$ 用来区分不同的场分量。例如，实标量场只有一个分量，复标量场可看作两个实场分量，四矢量场则带有时空指标。

时空坐标写为
\[
x^{\mu}=(ct,\bm{x}).
\]
在通常的场变分中，时空坐标 $x^{\mu}$ 作为自变量保持固定，真正被变分的是场函数 $\phi^{a}(x)$。因此，从粒子力学到场论的核心变化可以概括为
\[
q^{i}(t)
\quad\longrightarrow\quad
\phi^{a}(x^{\mu}).
\]
前者描述有限自由度系统，后者描述无限自由度系统。

粒子力学中的基本对象是拉格朗日量 $L$。场论中更自然的基本对象是拉格朗日密度
\[
\mathcal{L}=\mathcal{L}(\phi^{a},\partial_{\mu}\phi^{a},x^{\mu}).
\]
在某一惯性系中，总拉格朗日量由空间积分给出：
\[
L=\int \mathcal{L}\,\dif^{3}x.
\]
若采用四维体元
\[
\dif^{4}x=\dif x^{0}\,\dif x^{1}\,\dif x^{2}\,\dif x^{3}=c\,\dif t\,\dif^{3}x,
\]
则作用量写为式 \eqref{eq:ch7-action}：
\begin{equation}
	S=\frac{1}{c}\int_{\Omega}\mathcal{L}(\phi^{a},\partial_{\mu}\phi^{a},x^{\mu})\,\dif^{4}x.
	\label{eq:ch7-action}
\end{equation}
这里因子 $1/c$ 来自约定 $x^{0}=ct$。若采用 $x^{0}=t$，或把四维体元直接定义为 $\dif t\,\dif^{3}x$，则该因子会相应改变。整体常数不影响场方程，但在单位制上必须保持一致。

本章的目标是建立经典场论的三条基本线索：由作用量变分得到场方程；由局域连续性方程表达守恒律；由连续对称性得到 Noether 守恒流。

\section{场的作用量原理}

设场变量为 $\phi^{a}(x)$，作用量取为
\[
S=\frac{1}{c}\int_{\Omega}
\mathcal{L}(\phi^{a},\partial_{\mu}\phi^{a},x^{\mu})\,\dif^{4}x,
\]
其中 $\Omega$ 是四维时空中的积分区域。哈密顿原理要求真实场构型满足
\begin{equation}
	\delta S=0.
	\label{eq:ch7-stationary-action}
\end{equation}

对场作无穷小变分：
\[
\phi^{a}(x)\rightarrow \phi^{a}(x)+\delta\phi^{a}(x).
\]
在这种变分中，时空坐标固定不变。因此有
\[
\delta(\partial_{\mu}\phi^{a})=\partial_{\mu}(\delta\phi^{a}).
\]
同时要求场的变分在积分区域边界上消失：
\[
\delta\phi^{a}\big|_{\partial\Omega}=0.
\]

作用量的一阶变分为
\[
\delta S
=
\frac{1}{c}
\int_{\Omega}
\left[
\frac{\partial\mathcal{L}}{\partial\phi^{a}}\delta\phi^{a}
+
\frac{\partial\mathcal{L}}{\partial(\partial_{\mu}\phi^{a})}
\delta(\partial_{\mu}\phi^{a})
\right]\dif^{4}x.
\]
利用 $\delta(\partial_{\mu}\phi^{a})=\partial_{\mu}(\delta\phi^{a})$，得到
\[
\delta S
=
\frac{1}{c}
\int_{\Omega}
\left[
\frac{\partial\mathcal{L}}{\partial\phi^{a}}\delta\phi^{a}
+
\frac{\partial\mathcal{L}}{\partial(\partial_{\mu}\phi^{a})}
\partial_{\mu}(\delta\phi^{a})
\right]\dif^{4}x.
\]

第二项可用恒等式
\[
\frac{\partial\mathcal{L}}{\partial(\partial_{\mu}\phi^{a})}
\partial_{\mu}(\delta\phi^{a})
=
\partial_{\mu}\left[
\frac{\partial\mathcal{L}}{\partial(\partial_{\mu}\phi^{a})}
\delta\phi^{a}
\right]
-
\partial_{\mu}\left[
\frac{\partial\mathcal{L}}{\partial(\partial_{\mu}\phi^{a})}
\right]\delta\phi^{a}
\]
作分部积分。代入上式后得到
\[
\delta S
=
\frac{1}{c}
\int_{\Omega}
\left[
\frac{\partial\mathcal{L}}{\partial\phi^{a}}
-
\partial_{\mu}\left(
\frac{\partial\mathcal{L}}{\partial(\partial_{\mu}\phi^{a})}
\right)
\right]\delta\phi^{a}\,\dif^{4}x
+
\frac{1}{c}
\int_{\Omega}
\partial_{\mu}\left[
\frac{\partial\mathcal{L}}{\partial(\partial_{\mu}\phi^{a})}
\delta\phi^{a}
\right]\dif^{4}x.
\]
第二项是四维全散度项。由四维高斯定理，
\[
\int_{\Omega}
\partial_{\mu}\left[
\frac{\partial\mathcal{L}}{\partial(\partial_{\mu}\phi^{a})}
\delta\phi^{a}
\right]\dif^{4}x
=
\oint_{\partial\Omega}
\frac{\partial\mathcal{L}}{\partial(\partial_{\mu}\phi^{a})}
\delta\phi^{a}\,\dif\Sigma_{\mu}.
\]
由于边界上 $\delta\phi^{a}=0$，边界项消失。因此
\[
\delta S
=
\frac{1}{c}
\int_{\Omega}
\left[
\frac{\partial\mathcal{L}}{\partial\phi^{a}}
-
\partial_{\mu}\left(
\frac{\partial\mathcal{L}}{\partial(\partial_{\mu}\phi^{a})}
\right)
\right]\delta\phi^{a}\,\dif^{4}x.
\]
由于 $\delta\phi^{a}$ 在区域内部任意，得到场的 Euler--Lagrange 方程：
\begin{equation}
	\partial_{\mu}\left(
	\frac{\partial\mathcal{L}}{\partial(\partial_{\mu}\phi^{a})}
	\right)
	-
	\frac{\partial\mathcal{L}}{\partial\phi^{a}}
	=0.
	\label{eq:ch7-field-el}
\end{equation}
式 \eqref{eq:ch7-field-el} 就是经典场论中的基本变分方程。给定拉格朗日密度 $\mathcal{L}$ 后，场方程由这条公式确定。

\section{标量场的例子}

先考虑实标量场
\[
\phi=\phi(x^{\mu}).
\]
自由无质量标量场的拉格朗日密度可取为
\begin{equation}
	\mathcal{L}
	=
	\frac{1}{2}\partial_{\alpha}\phi\,\partial^{\alpha}\phi.
	\label{eq:ch7-massless-scalar-lagrangian}
\end{equation}
它不显含 $\phi$，所以
\[
\frac{\partial\mathcal{L}}{\partial\phi}=0.
\]
另一方面，由
\[
\mathcal{L}
=
\frac{1}{2}\eta^{\alpha\beta}(\partial_{\alpha}\phi)(\partial_{\beta}\phi)
\]
可得
\[
\frac{\partial\mathcal{L}}{\partial(\partial_{\mu}\phi)}=\partial^{\mu}\phi.
\]
代入式 \eqref{eq:ch7-field-el}，得到
\[
\partial_{\mu}\partial^{\mu}\phi=0.
\]
记达朗贝尔算子为
\begin{equation}
	\dalembert=\partial_{\mu}\partial^{\mu}.
	\label{eq:ch7-dalembert-def}
\end{equation}
在号差约定
\[
\eta_{\mu\nu}=\mathrm{diag}(1,-1,-1,-1)
\]
下，有
\begin{equation}
	\dalembert
	=
	\frac{1}{c^{2}}\frac{\partial^{2}}{\partial t^{2}}
	-
	\bm{\nabla}^{2}.
	\label{eq:ch7-dalembert-expanded}
\end{equation}
因此无质量自由标量场满足
\begin{equation}
	\dalembert\phi=0.
	\label{eq:ch7-massless-wave-equation}
\end{equation}
展开后即
\[
\left(
\frac{1}{c^{2}}\frac{\partial^{2}}{\partial t^{2}}
-
\bm{\nabla}^{2}
\right)\phi=0.
\]
这是以光速传播的自由波动方程。

若考虑有质量实标量场，为避免过早引入量子常数，可先把质量参数写成逆长度量纲的常数 $\mu$，取
\begin{equation}
	\mathcal{L}
	=
	\frac{1}{2}\partial_{\alpha}\phi\,\partial^{\alpha}\phi
	-
	\frac{1}{2}\mu^{2}\phi^{2}.
	\label{eq:ch7-massive-scalar-lagrangian}
\end{equation}
此时
\[
\frac{\partial\mathcal{L}}{\partial\phi}
=
-\mu^{2}\phi,
\qquad
\frac{\partial\mathcal{L}}{\partial(\partial_{\mu}\phi)}
=
\partial^{\mu}\phi.
\]
场方程为
\[
\partial_{\mu}\partial^{\mu}\phi+\mu^{2}\phi=0,
\]
即
\begin{equation}
	(\dalembert+\mu^{2})\phi=0.
	\label{eq:ch7-kg-mu}
\end{equation}
若进一步联系量子场论中粒子质量 $m$ 的记号，可取
\[
\mu=\frac{mc}{\hbar},
\]
则方程写为
\begin{equation}
	\left(\dalembert+\frac{m^{2}c^{2}}{\hbar^{2}}\right)\phi=0.
	\label{eq:ch7-kg-mass}
\end{equation}
这就是 Klein--Gordon 型方程。这里仍将它理解为经典场方程。

\section{多分量场与向量场的 Euler--Lagrange 方程}

一般场可以有多个分量，记为
\[
\phi^{a}(x).
\]
其中 $a$ 可以表示内部指标，也可以表示场分量标签。若拉格朗日密度为
\[
\mathcal{L}=\mathcal{L}(\phi^{a},\partial_{\mu}\phi^{a},x),
\]
则每一个场分量都满足与式 \eqref{eq:ch7-field-el} 同形的方程：
\begin{equation}
	\partial_{\mu}\left(
	\frac{\partial\mathcal{L}}{\partial(\partial_{\mu}\phi^{a})}
	\right)
	-
	\frac{\partial\mathcal{L}}{\partial\phi^{a}}
	=0.
	\label{eq:ch7-multifield-el}
\end{equation}

四矢量场可写为
\[
A_{\mu}=A_{\mu}(x).
\]
若
\[
\mathcal{L}=\mathcal{L}(A_{\mu},\partial_{\nu}A_{\mu},x),
\]
对 $A_{\mu}$ 作变分：
\[
A_{\mu}\rightarrow A_{\mu}+\delta A_{\mu}.
\]
重复前面的分部积分过程，可得
\begin{equation}
	\partial_{\nu}\left(
	\frac{\partial\mathcal{L}}{\partial(\partial_{\nu}A_{\mu})}
	\right)
	-
	\frac{\partial\mathcal{L}}{\partial A_{\mu}}
	=0.
	\label{eq:ch7-vector-field-el}
\end{equation}
这与标量场的公式完全同构，差别只在于场变量本身带有指标。后续电磁场理论中，四维势 $A_{\mu}$ 和场强张量都将使用这一类变分公式。

\section{拉格朗日密度的不唯一性}

粒子力学中，如果对拉格朗日量加上全时间导数
\[
L\rightarrow L+\frac{\dif F}{\dif t},
\]
则作用量只改变端点项，不改变运动方程。场论中有完全类似的结论。

若拉格朗日密度改变为
\[
\mathcal{L}\rightarrow \mathcal{L}+\partial_{\mu}K^{\mu},
\]
则作用量改变为
\[
S\rightarrow S+\frac{1}{c}\int_{\Omega}\partial_{\mu}K^{\mu}\,\dif^{4}x.
\]
由四维高斯定理，
\[
\int_{\Omega}\partial_{\mu}K^{\mu}\,\dif^{4}x
=
\oint_{\partial\Omega}K^{\mu}\,\dif\Sigma_{\mu}.
\]
这是边界项。在通常的变分问题中，场的边界值固定，或者要求场在无穷远处足够快衰减，因此这个边界项不影响 Euler--Lagrange 方程。

所以，拉格朗日密度具有式 \eqref{eq:ch7-lagrangian-density-equivalence} 所示的不唯一性：
\begin{equation}
	\mathcal{L}\sim \mathcal{L}+\partial_{\mu}K^{\mu}.
	\label{eq:ch7-lagrangian-density-equivalence}
\end{equation}
不同但只相差全散度项的拉格朗日密度给出相同的局域场方程。Noether 流和能动张量也可能因此相差一个改进项。

\section{局域守恒律与守恒荷}

对有限粒子体系，守恒律通常写成某个总量不随时间变化。对于连续场，守恒律更自然地写成局域形式。设某个物理量的密度为 $\rho$，通量为 $\bm{j}$。若该物理量没有源和汇，则对任意固定空间区域 $V$，有
\[
\frac{\dif}{\dif t}\int_{V}\rho\,\dif^{3}x
+
\oint_{\partial V}\bm{j}\cdot\dif\bm{S}=0.
\]
利用三维高斯定理，得到
\[
\int_{V}\left(
\frac{\partial\rho}{\partial t}
+\bm{\nabla}\cdot\bm{j}
\right)\dif^{3}x=0.
\]
由于区域 $V$ 任意，有
\begin{equation}
	\frac{\partial\rho}{\partial t}
	+\bm{\nabla}\cdot\bm{j}=0.
	\label{eq:ch7-continuity-3d}
\end{equation}
式 \eqref{eq:ch7-continuity-3d} 就是局域连续性方程。

采用 $x^{0}=ct$ 时，若定义四维守恒流为
\begin{equation}
	j^{\mu}=(c\rho,\bm{j}),
	\label{eq:ch7-four-current-convention}
\end{equation}
则
\[
\partial_{\mu}j^{\mu}
=
\frac{1}{c}\frac{\partial}{\partial t}(c\rho)
+\bm{\nabla}\cdot\bm{j}.
\]
因此式 \eqref{eq:ch7-continuity-3d} 可写为协变形式
\begin{equation}
	\partial_{\mu}j^{\mu}=0.
	\label{eq:ch7-local-conservation}
\end{equation}

在式 \eqref{eq:ch7-four-current-convention} 的约定下，与该守恒流对应的守恒荷应定义为
\begin{equation}
	Q=\int \rho\,\dif^{3}x
	=\frac{1}{c}\int j^{0}\,\dif^{3}x.
	\label{eq:ch7-conserved-charge}
\end{equation}
由 $\partial_{\mu}j^{\mu}=0$ 可得
\[
\frac{\partial j^{0}}{\partial t}=-c\,\partial_{i}j^{i}.
\]
于是
\[
\frac{\dif Q}{\dif t}
=
\frac{1}{c}\int \frac{\partial j^{0}}{\partial t}\,\dif^{3}x
=
-\int \partial_{i}j^{i}\,\dif^{3}x.
\]
若场在无穷远处足够快衰减，使边界通量为零，则
\[
\frac{\dif Q}{\dif t}=0.
\]
这说明四维局域守恒律 \eqref{eq:ch7-local-conservation} 蕴含式 \eqref{eq:ch7-conserved-charge} 中总荷 $Q$ 的守恒。

\section{场的 Noether 定理}

考虑一组场
\[
\phi^{a}(x).
\]
设发生无穷小变换
\[
x^{\mu}\rightarrow x'^{\mu}=x^{\mu}+\delta x^{\mu},
\]
\[
\phi^{a}(x)\rightarrow \phi'^{a}(x')=\phi^{a}(x)+\delta\phi^{a}.
\]
这里 $\delta\phi^{a}$ 表示总变化，即比较变换前点 $x$ 处的场值与变换后点 $x'$ 处的场值。

为了在同一个时空点比较场值，定义同点变化
\[
\bar{\delta}\phi^{a}=\phi'^{a}(x)-\phi^{a}(x).
\]
总变化与同点变化之间满足
\begin{equation}
	\bar{\delta}\phi^{a}
	=
	\delta\phi^{a}-\partial_{\nu}\phi^{a}\,\delta x^{\nu}.
	\label{eq:ch7-same-point-variation}
\end{equation}
本章采用如下约定：若无穷小变换由参数 $\epsilon^{A}$ 生成，则
\[
\delta x^{\mu}=\epsilon^{A}X_{A}^{\mu},
\qquad
\delta\phi^{a}=\epsilon^{A}\Delta_{A}\phi^{a}
\]
表示总变化，而 $\bar{\delta}\phi^{a}$ 表示同点变化。后文所有 Noether 流均按这一约定书写。

在场方程成立时，作用量的一阶变分可以化为边界项。对于严格对称变换，局域守恒律可写为
\[
\partial_{\mu}\left[
\mathcal{L}\,\delta x^{\mu}
+
\frac{\partial\mathcal{L}}{\partial(\partial_{\mu}\phi^{a})}
\bar{\delta}\phi^{a}
\right]=0.
\]
代入
\[
\bar{\delta}\phi^{a}
=
\delta\phi^{a}-\partial_{\nu}\phi^{a}\,\delta x^{\nu},
\]
得到
\[
\partial_{\mu}\left[
\mathcal{L}\,\delta x^{\mu}
+
\frac{\partial\mathcal{L}}{\partial(\partial_{\mu}\phi^{a})}
\delta\phi^{a}
-
\frac{\partial\mathcal{L}}{\partial(\partial_{\mu}\phi^{a})}
\partial_{\nu}\phi^{a}\,\delta x^{\nu}
\right]=0.
\]
将 $\delta x^{\mu}=\delta^{\mu}_{\nu}\delta x^{\nu}$ 写入，并定义典范能动张量
\begin{equation}
	T^{\mu}{}_{\nu}
	=
	\frac{\partial\mathcal{L}}{\partial(\partial_{\mu}\phi^{a})}
	\partial_{\nu}\phi^{a}
	-
	\delta^{\mu}_{\nu}\mathcal{L}.
	\label{eq:ch7-canonical-energy-momentum}
\end{equation}
利用式 \eqref{eq:ch7-canonical-energy-momentum}，守恒律可整理为
\[
\partial_{\mu}\left[
\frac{\partial\mathcal{L}}{\partial(\partial_{\mu}\phi^{a})}
\delta\phi^{a}
-
T^{\mu}{}_{\nu}\delta x^{\nu}
\right]=0.
\]

若写成由参数 $\epsilon^{A}$ 标记的守恒流，本章约定下可取
\begin{equation}
	(j^{\mu})_{A}
	=
	T^{\mu}{}_{\nu}X_{A}^{\nu}
	-
	\frac{\partial\mathcal{L}}{\partial(\partial_{\mu}\phi^{a})}
	\Delta_{A}\phi^{a}.
	\label{eq:ch7-noether-current}
\end{equation}
则有
\begin{equation}
	\partial_{\mu}(j^{\mu})_{A}=0.
	\label{eq:ch7-noether-conservation}
\end{equation}
式 \eqref{eq:ch7-noether-current} 与 \eqref{eq:ch7-noether-conservation} 就是场论中的 Noether 定理。

若拉格朗日密度在变换下不是严格不变，而是改变一个全散度，则守恒流需要相应加入边界修正项。实际使用时，最可靠的做法是从作用量变分的边界项直接读出守恒流。

拉格朗日形式的一个核心作用，是把对称性和守恒律联系起来。这个联系由 Noether 定理给出。其基本思想是：每一个连续对称性都对应一个守恒流。

这里有两个关键词。第一，对称性，指的是对系统变量作某种变换后，作用量不变，或者至多改变一个边界项。第二，连续性，指的是这种变换可以由无穷小参数连续生成。例如时间平移、空间平移、旋转、洛伦兹变换、全局相位变换都是连续变换。

离散变换，例如空间反演和时间反演，不属于 Noether 第一定理的直接讨论范围。在场论中，Noether 定理特别自然，因为守恒律自动以局域连续性方程的形式出现：
\[
\partial_{\mu}j^{\mu}=0.
\]
该式进一步给出守恒荷
\[
Q=\frac{1}{c}\int j^{0}\,\dif^{3}x.
\]

\section{能动张量}

\subsection{时空平移与典范能动张量}

考虑时空平移
\[
x^{\mu}\rightarrow x'^{\mu}=x^{\mu}+\epsilon^{\mu}.
\]
在本章采用的总变化约定下，平移可取
\[
\delta x^{\mu}=\epsilon^{\mu},
\qquad
\delta\phi^{a}=0.
\]
于是 Noether 流为
\[
(j^{\mu})_{\nu}=T^{\mu}{}_{\nu}.
\]
因此时空平移对称性给出守恒律
\begin{equation}
	\partial_{\mu}T^{\mu}{}_{\nu}=0.
	\label{eq:ch7-canonical-em-conservation}
\end{equation}
其中
\[
T^{\mu}{}_{\nu}
=
\frac{\partial\mathcal{L}}{\partial(\partial_{\mu}\phi^{a})}
\partial_{\nu}\phi^{a}
-
\delta^{\mu}_{\nu}\mathcal{L}
\]
即式 \eqref{eq:ch7-canonical-energy-momentum} 中的典范能动张量。指标 $\nu=0$ 对应能量守恒，$\nu=1,2,3$ 对应动量守恒。典范能动张量是时空平移对称性的 Noether 流。

\subsection{能动张量的不唯一性与改进}

由 Noether 定理直接得到的能动张量称为典范能动张量。前一节自然得到混合形式 $T^{\mu}{}_{\nu}$。为了讨论对称性和改进项，常将第二个指标升起，记
\begin{equation}
	T^{\mu\nu}=T^{\mu}{}_{\rho}\eta^{\rho\nu}.
	\label{eq:ch7-raise-em-index}
\end{equation}
将式 \eqref{eq:ch7-canonical-em-conservation} 升指标后，对应的守恒方程写为
\[
\partial_{\mu}T^{\mu\nu}=0.
\]

典范能动张量并不总是具有最理想的物理性质。它可能不是对称张量，也可能不具备某些规范不变性。这不表示 Noether 定理有问题，而是因为满足守恒方程的能动张量并不唯一。

若向 $T^{\mu\nu}$ 加上一个全散度项
\begin{equation}
	\Theta^{\mu\nu}=T^{\mu\nu}+\partial_{\rho}\chi^{\rho\mu\nu}.
	\label{eq:ch7-improved-em-tensor}
\end{equation}
并要求
\[
\chi^{\rho\mu\nu}=-\chi^{\mu\rho\nu},
\]
则
\[
\partial_{\mu}\partial_{\rho}\chi^{\rho\mu\nu}=0.
\]
这是因为偏导数对 $\mu,\rho$ 对称，而 $\chi^{\rho\mu\nu}$ 对 $\rho,\mu$ 反对称。因此
\[
\partial_{\mu}\Theta^{\mu\nu}
=
\partial_{\mu}T^{\mu\nu}.
\]
所以由式 \eqref{eq:ch7-improved-em-tensor} 定义的新能动张量仍满足同样的守恒方程：
\[
\partial_{\mu}\Theta^{\mu\nu}=0.
\]
若场在无穷远处衰减足够快，散度改进项对总能量和总动量的空间积分没有贡献。

这种自由度称为能动张量的改进。通过适当选择 $\chi^{\rho\mu\nu}$，可以把典范能动张量改写成更适合物理解释的形式。例如，对具有自旋的场，常用 Belinfante--Rosenfeld 方法得到对称能动张量。电磁场中也会出现类似问题：直接由 Noether 定理得到的典范能动张量可能显含势函数，而物理上更自然的能动张量应当用场强张量表示。

\section{洛伦兹对称性与角动量流}

时空平移给出能量动量守恒，洛伦兹对称性则给出角动量和 boost 相关守恒量。无穷小洛伦兹变换可写为
\[
\delta x^{\mu}=\omega^{\mu}{}_{\nu}x^{\nu},
\]
其中参数满足
\[
\omega_{\mu\nu}=-\omega_{\nu\mu}.
\]
对这种连续对称性应用 Noether 定理，可以得到广义角动量流
\[
M^{\rho\mu\nu}.
\]
它对后两个指标反对称：
\[
M^{\rho\mu\nu}=-M^{\rho\nu\mu}.
\]
守恒律为
\begin{equation}
	\partial_{\rho}M^{\rho\mu\nu}=0.
	\label{eq:ch7-angular-current-conservation}
\end{equation}

对于没有内禀自旋结构的标量场，角动量流只有式 \eqref{eq:ch7-orbital-angular-momentum-current} 所示的轨道部分：
\begin{equation}
	M^{\rho\mu\nu}
	=
	x^{\mu}T^{\rho\nu}-x^{\nu}T^{\rho\mu}.
	\label{eq:ch7-orbital-angular-momentum-current}
\end{equation}
这正是粒子力学中 $\bm{L}=\bm{r}\times\bm{p}$ 的四维场论推广。

对矢量场、旋量场等带有内禀洛伦兹变换结构的场，还会出现自旋部分：
\[
M^{\rho\mu\nu}
=
x^{\mu}T^{\rho\nu}
-
x^{\nu}T^{\rho\mu}
+
S^{\rho\mu\nu}.
\]
其中 $S^{\rho\mu\nu}$ 称为自旋角动量流。它来自场分量本身在洛伦兹变换下的变化，而不是来自坐标位置因子。

因此，场的总角动量一般由轨道角动量和内禀自旋角动量组成。若使用经过 Belinfante--Rosenfeld 对称化的能动张量，自旋部分可以被吸收到改进后的能动张量中，使总角动量流写成更紧凑的形式。

\section{内部对称性与守恒流}

Noether 定理不仅适用于时空变换，也适用于不改变时空坐标的内部变换。这类对称性称为内部对称性。

最简单的例子是复标量场的全局相位变换。设复场为
\[
\psi(x),
\]
其复共轭为
\[
\psi^{*}(x).
\]
若拉格朗日密度在全局相位变换下不变：
\[
\psi\rightarrow e^{\mathrm{i}\alpha}\psi,
\qquad
\psi^{*}\rightarrow e^{-\mathrm{i}\alpha}\psi^{*},
\]
则系统具有全局 $U(1)$ 对称性。

对无穷小参数 $\alpha$，有
\[
\delta\psi=\mathrm{i}\alpha\psi,
\qquad
\delta\psi^{*}=-\mathrm{i}\alpha\psi^{*}.
\]
这里时空坐标不变，$\delta x^{\mu}=0$。若把 $\psi$ 与 $\psi^{*}$ 看作独立场变量，则由式 \eqref{eq:ch7-noether-current} 可得 Noether 流
\begin{equation}
	j^{\mu}
	=
	\mathrm{i}
	\left[
	\frac{\partial\mathcal{L}}{\partial(\partial_{\mu}\psi)}\psi
	-
	\frac{\partial\mathcal{L}}{\partial(\partial_{\mu}\psi^{*})}\psi^{*}
	\right].
	\label{eq:ch7-u1-current}
\end{equation}
它满足
\[
\partial_{\mu}j^{\mu}=0.
\]
该流的整体符号取决于全局相位变换采用 $\psi\rightarrow e^{\mathrm{i}\alpha}\psi$ 还是 $\psi\rightarrow e^{-\mathrm{i}\alpha}\psi$。物理电荷的正负号也需要与这一约定一致。

对应的守恒荷仍按式 \eqref{eq:ch7-conserved-charge} 定义。
若场在空间无穷远处衰减足够快，则 $\dif Q/\dif t=0$。

这里需要区分两类对称性。全局对称性的变换参数是常数，例如 $\alpha=\mathrm{const}$；局域对称性的变换参数可以依赖时空，例如 $\alpha=\alpha(x)$。全局对称性通过 Noether 定理给出守恒流。局域对称性则进一步要求引入规范场，并导致规范相互作用。电磁场正是 $U(1)$ 规范对称性的典型例子。

\section{本章小结}

本章建立了经典场论的拉格朗日形式。粒子力学中的有限自由度 $q^{i}(t)$ 被推广为时空中的场变量 $\phi^{a}(x)$，拉格朗日量被推广为拉格朗日密度 $\mathcal{L}$。在 $x^{0}=ct$ 的约定下，作用量写为
\[
S=\frac{1}{c}\int \mathcal{L}\,\dif^{4}x.
\]

作用量驻值条件给出场的 Euler--Lagrange 方程：
\[
\partial_{\mu}\left(
\frac{\partial\mathcal{L}}{\partial(\partial_{\mu}\phi^{a})}
\right)
-
\frac{\partial\mathcal{L}}{\partial\phi^{a}}
=0.
\]
例如，实标量场的自由拉格朗日密度给出无质量波动方程或有质量 Klein--Gordon 型方程。

场论中的局域守恒律写为
\[
\partial_{\mu}j^{\mu}=0.
\]
若 $j^{\mu}=(c\rho,\bm{j})$，则对应守恒荷为
\[
Q=\frac{1}{c}\int j^{0}\,\dif^{3}x.
\]
在边界通量为零时，$Q$ 不随时间变化。

Noether 定理说明，连续对称性给出守恒流。时空平移对称性给出典范能动张量
\[
T^{\mu}{}_{\nu}
=
\frac{\partial\mathcal{L}}{\partial(\partial_{\mu}\phi^{a})}
\partial_{\nu}\phi^{a}
-
\delta^{\mu}_{\nu}\mathcal{L},
\]
并满足
\[
\partial_{\mu}T^{\mu}{}_{\nu}=0.
\]
洛伦兹对称性给出角动量流，内部连续对称性给出相应的内部荷流。例如，复标量场的全局 $U(1)$ 相位对称性给出守恒流。

本章的核心思想可以概括为
\[
\delta S=0
\quad\Longrightarrow\quad
\text{场方程},
\]
以及
\[
\text{连续对称性}
\quad\Longrightarrow\quad
\text{守恒流}.
\]
后续电磁场理论将以本章为基础：电磁场的拉格朗日密度给出麦克斯韦方程，时空平移对称性给出电磁场能动张量，规范对称性则与电荷守恒密切相关。
